% TOP_PROBLEM: {"id": 2160, "title": "Erdős Problem #529", "category_id": 3, "category": {"id": 3, "name": "graph_theory", "display_name": "Graph Theory", "description": "Problems involving graphs, networks, and their properties.", "slug": "graph-theory", "order_index": 3, "created_at": "2026-07-31T15:26:25.671Z"}, "status": "open"}
% FIRST_POSTED: null
% ENTRY_KIND: statement_only
% Imported from: unsolvedmath_additional_500.tex; UnsolvedMath ID 2160
% !TeX program = lualatex
% Compile twice: lualatex 2160.tex
% Self-contained: no external bibliography, images, or \input files.
% Numeric section labels and visible numbers are original UnsolvedMath IDs.
\documentclass[11pt,a4paper]{article}
\usepackage[margin=24mm,headheight=14pt]{geometry}
\usepackage{fontspec}
\setmainfont{Latin Modern Roman}
\setsansfont{Latin Modern Sans}
\setmonofont{Latin Modern Mono}
\usepackage{amsmath,amssymb,mathtools}
\usepackage{unicode-math}
\setmathfont{Latin Modern Math}
\usepackage{microtype}
\usepackage{enumitem,xurl,needspace}
\usepackage[dvipsnames]{xcolor}
\usepackage{hyperref}
\hypersetup{unicode=true,colorlinks=true,linkcolor=MidnightBlue,urlcolor=MidnightBlue,
 pdftitle={UnsolvedMath Problem 2160},pdfauthor={Source-annotated compilation},bookmarksnumbered=true}
\usepackage{bookmark,tocloft,titlesec,fancyhdr}
\setlength{\cftsecnumwidth}{4.2em}
\cftsetpnumwidth{2.5em}
\cftsetrmarg{4em}
\titleformat{\section}{\Large\bfseries\raggedright}{\thesection}{0.7em}{}
\pagestyle{fancy}\fancyhf{}
\fancyhead[L]{\small\sffamily Additional UnsolvedMath statements}
\fancyhead[R]{\small Original catalogue IDs}
\fancyfoot[C]{\thepage}
\setlength{\parindent}{0pt}
\setlength{\parskip}{5pt plus 1pt}
\setlength{\emergencystretch}{3em}
\setcounter{tocdepth}{1}\setcounter{secnumdepth}{1}
\setlist[itemize]{leftmargin=3em,itemsep=4pt,topsep=4pt,parsep=0pt}
\allowdisplaybreaks
\newcommand{\N}{\mathbb{N}}
\newcommand{\Z}{\mathbb{Z}}
\newcommand{\Q}{\mathbb{Q}}
\newcommand{\R}{\mathbb{R}}
\newcommand{\C}{\mathbb{C}}
\newcommand{\F}{\mathbb{F}}
\newcommand{\A}{\mathbb{A}}
\newcommand{\PP}{\mathbb{P}}
\newcommand{\E}{\mathbb{E}}
\newcommand{\eps}{\varepsilon}
\newcommand{\dd}{\,\mathrm d}
\newcommand{\sref}[2]{\hyperlink{source:#1:#2}{[S#2]}}
\newif\ifshowcatalogue
\showcataloguetrue
\newcommand{\cataloguescope}[1]{\ifshowcatalogue
 \paragraph{Statement recorded in the supplied source.}
 #1\fi}
\begin{document}
% UnsolvedMath ID 2160; source catalogue code EP-529

\setcounter{section}{2159}

\section{Expected Displacement of Self-Avoiding Walks}\label{2160}

{\small\sffamily UnsolvedMath ID 2160 · Graph Theory}

\subsection{Definitions and mathematical statement}

\paragraph{Shared notation.}Unless a section says otherwise, $\N=\{1,2,\ldots\}$,
$\Z$ is the ring of integers, $\Q,\R,\C$ are the rational, real, and complex
fields, $[N]=\{1,\ldots,\lfloor N\rfloor\}$, and $\log$ is the natural
logarithm. For real functions of a parameter tending to infinity,
$f\ll g$ and $f=O(g)$ mean $|f|\leq Cg$ eventually for some fixed $C$;
$f\gg g$ reverses this comparison; $f=o(g)$ means $f/g\to0$;
$f\sim g$ means $f/g\to1$; and $f\asymp g$ means both order bounds.
A subscript on $O$ or $\ll$ specifies allowed constant dependence.
The expression $n^{o(1)}$ in an upper bound means growth smaller than
$n^\epsilon$ for each fixed $\epsilon>0$. Local definitions govern any
exception, finite-field convention, or use of the same symbol for another
object. An asymptotic claim is not automatically an effective algorithm.



Choose uniformly among all nearest-neighbour self-avoiding walks of exactly $n$ steps from the origin. Equivalently, condition simple random walk of that length on no repeated vertex. The observable is the expected Euclidean norm of the endpoint, not its squared norm. \sref{2160}{1} \sref{2160}{2}

\paragraph{Statement recorded in the supplied source.}
Let $d_k(n)$ be the expected distance from the origin after taking $n$ random steps from the origin in $\mathbb{Z}^k$ (conditional on no self intersections) - that is, a self-avoiding walk. Is it true that $ \lim_{n\to \infty}\frac{d_2(n)}{n^{1/2}}= \infty? $ Is it true that $ d_k(n)\ll n^{1/2} $ for $k\geq 3$? \sref{2160}{1}

\subsection{Short English statement}

How far does a uniformly sampled self-avoiding lattice walk typically end from its starting point, compared with the square root of its length?

\subsection{Sources}

{\small

\begin{itemize}

\item[\hypertarget{source:2160:1}{S1}] ULAM.AI, \emph{UnsolvedMath}, supplied \texttt{problems.json}, record 2160 (EP-529), statement and background. The embedded literature-review date is 2026-08-17; that is the archive’s review date, not a fresh certification. Dataset landing page: \url{https://huggingface.co/datasets/ulamai/UnsolvedMath}.

\item[\hypertarget{source:2160:2}{S2}] T. F. Bloom, \emph{Erdős Problems}, Problem 529, \url{https://www.erdosproblems.com/529}. Reference imported from the supplied record; the live problem page was not independently checked for this compilation.

\item[\hypertarget{source:2160:3}{S3}] Takashi Hara and Gordon Slade, Self-avoiding walk in five or more dimensions. I. The critical behaviour, Communications in Mathematical Physics 147 (1992), 101-136. \url{https://doi.org/10.1007/BF02099530}. Bibliographic information imported from the supplied record.

\item[\hypertarget{source:2160:4}{S4}] Hugo Duminil-Copin and Alan Hammond, Self-avoiding walk is sub-ballistic, arXiv:1205.0401 (2012). \url{https://arxiv.org/abs/1205.0401}. Bibliographic information imported from the supplied record.

\item[\hypertarget{source:2160:5}{S5}] Duminil-Copin, Hugo and Hammond, Alan, Self-avoiding walk is sub-ballistic. Comm. Math. Phys. (2013), 401--423. Bibliographic information imported from the supplied record.

\item[\hypertarget{source:2160:6}{S6}] Hara, Takashi and Slade, Gordon, Critical behaviour of self-avoiding walk in five or more dimensions. Bull. Amer. Math. Soc. (N.S.) (1991), 417--423. Bibliographic information imported from the supplied record.

\end{itemize}

}

\label{end:2160}
\end{document}
